\section{The real \texorpdfstring{$S_3$}{S3} anchor and the detector at infinity}\label{sec:real}

Suppose that \(W=E[2]\) has image \(S_3\) and that its quadratic
sign field \(L\) is real. Complex conjugation then acts trivially on
\(W\): its image has order at most two, and the only even permutation
of that order in \(S_3\) is the identity. We prove the remaining
\(S_3\) case of Proposition~\ref{prop:anchor}.

\begin{proposition}[The real \(S_3\) anchor]\label{real:anchor}
There are a fundamental discriminant \(h\) prime to \(2N_E\), allowing
\(h=1\), and an imaginary quadratic field \(K\) of discriminant \(-d\)
prime to \(2hN_E\), in which every prime dividing \(2hN_E\) splits, such
that
\[
 \an(E^h)+\an(E^{h(-d)})=1,
 \qquad X(E^h)+X(E^{h(-d)})=0.
\]
\end{proposition}

The new ingredient is a nonzero residual functional on the real
components of a Fricke twist of a modular Jacobian. To use that
functional, we construct an integral global class by a Pfaffian
cofactor calculation. The horizontal logarithm calculation and the
integral denominator calculation are kept separate until the final
specialization.

\subsection{The level and an elliptic starting pair}

The representation \(W\) is absolutely simple. Over \(\mathbb F_4\) it
is induced from a cubic character \(\psi\) of \(G_L\), and the
nontrivial automorphism \(\iota\) of \(L/\mathbb Q\) inverts \(\psi\).
Write \(d_L>0\) for the fundamental discriminant of \(L\).

\begin{lemma}[A level supporting the real character]\label{real:level}
The character \(\psi\) is a character of an ordinary ring class group
of \(L\). If \(b\) is its ring conductor, then its character conductor
ideal is \(b\mathcal O_L\), and the odd part of the Artin conductor of
\(W\) is \(d_L^{\rm odd}(b^{\rm odd})^2\).
Choose a good prime \(\ell_0\equiv3\pmod4\), outside the conductor
support, whose Frobenius on \(W\) is a \(3\)-cycle. Then
\begin{equation}\label{real:level-formula}
 n=2^\delta d_L^{\rm odd}(b^{\rm odd}\ell_0)^2,
 \qquad \delta=v_2(d_L)\bmod2
\end{equation}
satisfies \(\mathbb Q(\sqrt n)=L\), and \(\psi\) factors through
\(\Pic(\mathbb Z[\sqrt n])\).
\end{lemma}

\begin{proof}
The character has no conductor at infinity, since its order is three.
It kills rational ideles: anti-invariance makes its value on a rational
idele equal to its inverse, and an element of order dividing both two
and three is trivial. At split primes the two conductor exponents
agree. At a ramified odd prime a positive conductor exponent is even.
Indeed conjugation acts trivially on residue units and by \((-1)^j\)
on the last contributing principal-unit layer of index \(j\);
anti-invariance forces \(j\) to be odd, so the conductor exponent
\(j+1\) is even. If \(2\) ramifies in \(L\), wild inertia in the
\(S_3\) image contains a transposition. Its normalizer has order two,
so the decomposition group has order two and the cubic character is
unramified there. If \(2\) is unramified in \(L\), the character
exponents at two are at most one, by tameness.

These local statements give conductor \(b\mathcal O_L\). Triviality
on the local units at the corresponding ray depth, together with
triviality on rational units, implies triviality on the units of the
order of conductor \(b\). Thus \(\psi\) factors through its ordinary
ring class group. Lift \(\psi\) to characteristic zero and use the
conductor formula for induction. The inertia-invariant dimensions
agree after reduction, including for an order-two subgroup; the odd
wild-invariant dimensions agree as well. This gives the asserted odd
Artin conductor.

Chebotarev supplies \(\ell_0\): the prescribed \(3\)-cycle and the
condition \(\ell_0\equiv3\pmod4\) are compatible because \(L\) is
real. Formula~\eqref{real:level-formula} preserves the squareclass of
\(d_L\). The conductor of \(\mathbb Z[\sqrt n]\) contains all required
odd factors and also contains two when \(d_L\) is odd. These are
exactly the remaining possible character-conductor factors.
\end{proof}

Put \(X=X_0(n)\), \(J=J_0(n)\), and let \(w_n\) be the Fricke
involution.

\begin{lemma}[A pair with Selmer corank one]\label{real:selmer-pair}
There is an odd fundamental discriminant \(h\), prime to \(2nN_E\)
and allowing \(h=1\), such that \(E_0=E^h\) has trivial two-Selmer
group. There is also a fresh prime \(d\equiv7\pmod8\) such that
\(K=\mathbb Q(\sqrt{-d})\) splits every prime dividing \(2nhN_E\),
\(W(\mathrm{Fr}_d)\) is a \(3\)-cycle, and
\[
 s_2(E_0/K)=1.
\]
The class number of \(K\) is odd and its units are precisely the signs.
\end{lemma}

\begin{proof}
First take an odd fundamental signed twist, prime to \(2nN_E\), with
functional sign plus; use a split negative discriminant if the sign
must change. Since \(W\) has no rational invariants, the two-Selmer
dimension is even by the parity comparison of
Section~\ref{sec:classical}.
While this dimension is at least two, take a positive good prime
\(p\equiv1\pmod4\), locally square at all earlier support, with trivial
Frobenius on \(W\), at which two independent old classes have rank-two
joint finite evaluation. Such a prime exists: restriction to the
kernel of \(W\) and of the abelian congruence data detects the classes
by the normal \(C_3\) without invariants. The joint evaluation image is
additive over \(\mathbb F_2\), and absolute simplicity and independence
make it simultaneously full. Chebotarev then imposes the required
evaluation and congruences.

At the new prime the old unramified plane and the new twist Kummer
plane are transverse. The Kummers of the two-torsion points give the
two inertia coordinates: untwisted halves are unramified, whereas
twist inertia negates them. Reciprocity makes the new ramification
images annihilate the old evaluations. The remaining kernel is the
common strict Selmer group, so the dimension drops by two. Iterating
produces \(E_0=E^h\) with trivial two-Selmer group.

Choose \(d\) by Chebotarev and quadratic reciprocity. The residue
conditions expressing splitting of all primes of \(2nhN_E\) in
\(K\) force \(d\) to split in \(L\), which is compatible with a
\(3\)-cycle on the abelian intersection. Genus theory gives odd class
number, and \(d\equiv7\pmod8\) excludes extra units. The local
conditions for \(E_0\) and \(E_0^{-d}\) agree at every finite place:
at the old support this follows from splitting, and at \(d\) the
local cohomology vanishes by the cycle Frobenius. At infinity one
Kummer line can change. The twist's two-Selmer dimension is therefore
at most one; the coprime split twist sign and parity make it exactly
one. Quadratic induction now gives \(s_2(E_0/K)=1\).
\end{proof}

\subsection{A Hecke eigenfunctional on real components}

Use a maximal-order split-level Heegner point of this \(K\) on \(X\) as in \eqref{eq:GZ}. For its single Hilbert trace subtract the trace of the cusp at infinity, writing \(Y\in J(K)\) for the result. Put \(Y'=(T_\ell-\ell-1)Y\) for another good prime outside \(2nd\) with cycle Frobenius. It is represented by applying that difference to the trace alone, with no cusps. This is a \(\sigma=w_n c\)-real divisor, \(c\) complex conjugation: both operations reverse orientations and permute the ideal labels into the reversed orbit (the Fricke quotient translates by a level ideal).

\begin{proposition}[The detector at infinity]\label{real:detector}
There is an \(\mathbb F_2\)-linear functional
\[
 \Theta:H^1(\langle\sigma\rangle,T_2J)\to\mathbb F_2,\quad
 \Theta T_p=\operatorname{tr}W(\mathrm{Fr}_p)\Theta\ (p\nmid2n),
 \qquad \Theta([Y'])\ne0 .                                             \tag{R0}\label{eq:R0}
\]
where \([Y']\) denotes the twisted real Kummer class.
\end{proposition}

\begin{proof}
The Betti exponential sequence and its 2-adic period lattice identify
this group with the components of \(J(\mathbb C)^\sigma\), or its
quotient by norms, compatibly with Kummer. We will assign a weight to
each real circle of \(X\), show that the resulting divisor count is
a Hecke eigenfunctional, and evaluate it on the Hilbert trace.

\emph{Weights on the real circles.}

A \(\sigma\)-fixed interior point gives an anti-linear isogeny
\(S\) of the elliptic torus, of degree \(n\) and with the specified
cyclic kernel: compose the quotient map with the anti-complex Fricke
identification. Since \(S^2=n\), the rational period lattice is a
line over \(L\), with \(\sqrt n\) acting by \(S\). Identifying
this line with \(L\) gives a fractional ideal \(I\subset L\),
well-defined up to \(L^\times\), with multiplier order
\(O\supset O_0:=\mathbb Z[\sqrt n]\).

The ideal \(I\) is locally invertible over \(O\). Indeed, locally
scale a quadratic-order lattice into the maximal order, with both
minimal valuations zero in the split case. It can then be scaled by
a unit to contain 1, avoiding at most two proper residual subspaces.
A lattice between \(\mathbb Z_p\) and the maximal order that
contains 1 is an order. This gives the asserted invertibility over
the multiplier order, and hence a class \(a=[I]\in\Pic(O)\).
No odd prime \(p\) divides \([O:O_0]\): otherwise the form of
nested quadratic orders gives \(p\mid n\) and makes \(S\)
scalar modulo \(p\), hence zero, contradicting its cyclic kernel.

Use the extension map \(\Pic(O_0)\to\Pic(O)\) to test whether
\(\psi\) descends to \(O\). Give the point weight zero if it
does not; if it does, continue to denote the descended character by
\(\psi\) and give the point weight
\[
                           \theta=\psi(a)+\psi(a)^{-1}
\]
in characteristic two. Replacing \(S\) by \(-S\) conjugates the
ideal class and therefore preserves this weight. There is no
order-two elliptic stabilizer modulo signs: multiplication by \(i\)
cannot preserve the level at \(\ell_0\). Multiplying \(S\) by a
possible order-three elliptic automorphism, up to sign, conjugates
it by a level automorphism up to sign, since the anti-complex
isogeny inverts these scalars. Thus these choices also preserve the
weight.

The weights are constant on each real circle. On a compact local
uniformizing chart, trivialize the integral period lattice. The
anti-linear maps of fixed degree have bounded norm, so only finitely
many integral matrices occur. Each branch therefore has constant
ideal-class data; at a branch meeting the possible choices have the
same weight by the preceding sign and automorphism calculation.
No cusp is fixed: its denominator divisor of \(n\) is unchanged
by \(c\) and exchanged with its Fricke complement by \(w_n\),
whereas \(n\) is nonsquare.

\emph{From circle weights to a component functional.}
There exists a real point: on the chosen Hilbert orbit \(\sigma\)
acts on the conductor-one labels by \(z\mapsto a_0z^{-1}\) for a
fixed \(a_0\). The class number is odd, so \(z^2=a_0\) has exactly
one solution. Denote the corresponding fixed CM point by \(x\).
Every invariant line bundle class therefore has a real divisor:
use the fiber at \(x\) to lift the real action with square one,
then apply real Riemann--Roch after adding a sufficiently large
multiple of \(x\).

For such a divisor, count its points on each real circle modulo two,
multiply by that circle's weight \(\theta\), and add. A principal
invariant divisor is the divisor of a function that can be scaled to
be real; its parity on each circle is even. Norm divisors also have
zero count. The weighted count consequently descends to degree-zero
divisor classes and vanishes on the identity component, which
consists of norms. This defines the required component functional.

\emph{The Hecke eigenlaw and the Hilbert trace.}
We can compute the Hecke action at generic real points, since small
motions preserve the component counts both before and after the
correspondence. Nonreal pairs contribute zero. On the unbranched
locus without automorphisms modulo signs, real lifts under \(T_p\)
correspond to \(S\)-stable \(p\)-lines. At inert \(p\) there are
none. At split \(p\), the descended anti-isogenies give the two
prime-ideal translates of \(I\), with the same multiplier order.
If \(\psi\) descends to this order and \(\mathfrak p\mid p\),
their weights add to
\[
 \bigl(\psi(\mathfrak p)+\psi(\mathfrak p)^{-1}\bigr)
 \bigl(\psi(a)+\psi(a)^{-1}\bigr).
\]
If \(\psi\) does not descend, all these weights are zero. This
proves \(\Theta T_p=\operatorname{tr}W(\mathrm{Fr}_p)\Theta\).

It remains to compute the weight of \(x\). The character \(\psi\)
factors through its multiplier order \(O\). The only case not
already settled by the odd-index argument is nonzero character
conductor at two when \(n\) and \(d_L\) are odd. Here \(S\)
interchanges the two parts of the split
\(\mathcal O_K\otimes\mathbb Z_2\)-lattice invertibly, by
anti-linear scalar conjugation. It cannot be scalar modulo two,
so the overorder index is odd in this case as well.

The CM action of \(\sqrt{-d}\) anticommutes with \(S\). On \(L\)
it therefore has the form \(\beta\iota\). Its image
\(\beta I^\iota\subset I\) is an \(O\)-submodule of index \(d\).
The prime \(d\) splits in \(L\) and is prime to the order
conductor, so this image is \(\mathfrak d I\) for one of the
primes \(\mathfrak d\) above \(d\). Taking ideal classes gives
\[
 [\mathfrak d]a=[I^\iota]=a^{-1},\qquad
                         [\mathfrak d]=a^{-2}.
\]
The choice of \(d\) gives \(\psi(\mathfrak d)\ne1\), whence
\(\theta(x)=1\). All other points in the Hilbert orbit occur in
nonreal pairs. Since the detector prime \(\ell\) has cycle
Frobenius and \(\ell+1\) is even, applying
\(T_\ell-\ell-1\) leaves this count equal to one. This proves
\eqref{eq:R0}.
\end{proof}

\subsection{The integral Hecke module}

Use the anemic cusp Hecke algebra at this level (all \(T_p,\ p\nmid2n\)). By \eqref{eq:R0} it has the maximal ideal of \(W\); write \(A\) for its completed factor and \(M=(T_2J)_A\). It is a reduced order finite free over \(\mathbb Z_2\) with residue \(\mathbb F_2\). Here and below we use ordinary newform theory and the modular abelian factors with their Tate realizations and level/conductor compatibility. 
\begin{lemma}[The rank-two Hecke module]\label{real:packet}
There is a free rank-two \(A\)-module \(P\), with Galois action of
cyclotomic determinant and residual representation \(W\), such that
\begin{equation}\label{real:morita}
 M=P\otimes_A D.
\end{equation}
Here Galois acts only on \(P\), \(D\) is a multiplicity lattice free
over \(\mathbb Z_2\), and \(P\) rationally carries the primitive
newform representations in the Hecke factor. No freeness of \(D\)
over \(A\) is asserted or needed.
All these primitive forms are non-CM. Their levels equal the specified
level away from \(2,\ell_0\); at \(2\mid n\) they are Steinberg.
The only possible multiplicity increase comes from
\(\ell_0^2\)-raising a form unramified at \(\ell_0\).
\end{lemma}

\begin{proof}
Rational plane traces of group elements belong to \(A\) reducing to traces of \(W\), by density. Componentwise reductions have semisimplification \(W\) after residue extension by trace and determinant, hence are absolutely irreducible. Choose four group elements spanning the residual matrix algebra (by absolute simplicity); their trace Gram determinant is a unit. They therefore span all the Galois matrices integrally, a free rank-four algebra acting on \(M\), reducing to the matrix algebra. Lift a primitive residual matrix idempotent by completeness. Its associated left ideal is free rank two and the algebra acts as all matrices on it by Nakayama. This gives the assertions including the decomposition (matrix Morita equivalence); \(D\) is free over \(\mathbb Z_2\), not required free over \(A\).

All occurring primitive forms are non-CM (else on an imaginary quadratic field the Tate reductions semisimplify reducibly, impossible for \(W\)). Their levels away from \(2,\ell_0\) are exactly the given one since residual conductor is a lower bound (inertia invariants, and odd wild invariants). At \(\ell_0\) exponent one is impossible by cycle Frobenius versus Steinberg with trivial nebentypus. Also at \(2\mid n\) no old form of good level occurs. Otherwise its modular abelian variety of good reduction at two would yield a finite flat model of the local \(W\): it occurs as a global two-torsion subquotient, by the rational Tate realizations and lattice independence of reduction semisimplification (the chosen coefficient embedding reduces to a scalar extension of \(W\)). Take closures and quotients over the DVR. The unique local line would give a finite flat subgroup of order two, thus \(\mu_2\) or \(\mathbb Z/2\) over \(\mathbb Z_2\) (order-two finite flat classification). The nontrivial generic point of the quotient extends by properness, so lifting it gives the residual quadratic radical by a torsor under that subgroup, a unit or unramified radical. This contradicts \(v_2(d_L)=3\) with decomposition of order two. Thus \(2\) is purely Steinberg then. The only multiplicity increase is possible by \(\ell_0^2\)-raising a primitive form unramified there.
\end{proof}

Fix a unimodular determinant form on \(P\). We will use the detector
by factoring the projected Hilbert trace as
\[
                         Y=a p,
\]
where \(p\) is an integral global class fixed by \(\sigma=w_nc\)
and \(a\) belongs to the coefficient order, after the flat extension
used for measures. If \(a\) has zero residue, applying
\(T_\ell-\ell-1\) and restricting to infinity contradicts
\eqref{eq:R0}. The construction below must therefore control the
integral class \(p\), not just its rational span. The Pfaffian
supplies a global lift; the subsequent logarithm and denominator
comparisons make its specialization integral.

\subsection{Parity and the Pfaffian strict system}

\begin{lemma}[Parity on the primitive components]\label{real:parity}
If \(g\) is one of the primitive components or its quadratic twist and has odd functional sign, its two-adic characteristic-zero Kummer Selmer dimension over \(\mathbb Q\) at the chosen embedding is odd.
\end{lemma}

\begin{proof}
Choose an odd imaginary discriminant \(k'\) splitting \(2N_gN_{E_0}\) with \(L(g\otimes\chi_{k'},1)\ne0\), by the split sign and twist nonvanishing as at the outset. This last form has zero Selmer rank there: use \eqref{eq:K0} with a further imaginary split Heegner companion giving simple twist zero by derivative nonvanishing. The Tate absolute irreducibility required holds over any such imaginary field (residual \(W\) with real quadratic \(L\)). This uses the split Gross--Zagier input to \eqref{eq:K0}. 

Over the field of \(k'\) use one full and one strict condition at each split pair over \(2N_gN_{E_0}\), comparing a Tate lattice for \(g\) and the elliptic \(T_2E_0\) at a common scalar DVR. Use \eqref{eq:H1} on induction including good unramified conditions at the odd discriminant places. These are perfect self-dual problems in degrees \(1,2\) (residual invariant vanishing, injectivity for conditions in degree zero and duality), with identical residual diagrams. Thus the characteristic-zero strict-pair dimensions have equal parities. Odd local factors are acyclic in characteristic zero by invariant vanishing for the abelian Tate realization and the odd Euler formula. At the split dyadic pair the full/strict plane and the plane of finite Kummer lines meet in a line, with no local invariants; both are Lagrangian in degree one for the conjugate pairing, which is symmetric (cup sign and alternating determinant form, with conjugate local invariants identified). For the form the finite lines are exact orthogonals by abelian polarization and Hecke projection, with Lie dimension one. Thus passing to the ordinary Selmer problem flips parity by the symmetric-plane comparison of the arithmetic diagram conventions. The ordinary dimension on \(E_0\) there is odd by quadratic induction, elliptic parity and split sign. Hence also on \(g\); decomposing by quadratic induction with the twist Selmer rank zero gives the claim.
\end{proof}

Return to \(K\) of discriminant \(-d\). Allow a fixed finite set \(S\) of split rational primes containing those of \(2nN_{E_0}\); the place over \(d\) when included for induction always has the original unramified condition. Use the split anticyclotomic free \(2\)-direction \(1+t\), with strict at \(w\mid2\) as before, and also product character variables \(1+u_j\) from the 2-Sylows at varying good inert prime conductors \(q_j\to-1\) two-adically. Add the conductor places to allowed support. By units and odd class number the relative cyclic Sylows have orders \(2^{v_2(q_j+1)}\), with primitive relative inertia actions, independent across the primes. All systems here use the finite models and common-stage extensions/limits of the arithmetic diagram conventions as in the other split construction. Put \(R=A[[t,\mathbf u]]\), with total scalar character \(\chi\) on the arithmetic coefficient \(P\chi\). The full strict determinant below will use strict at \(w\), full elsewhere over finite allowed places (not relaxing ramification at \(d\)). Define the modified complex \(C\) by the following local conditions.
\begin{itemize}
\item keep that dyadic pair; at each odd pair of \(S\) put one full coordinate (chosen with \(\chi(\mathrm{Fr}_v)=\xi_q\) equal to the positive CM orientation's inverse-translation weight as at \eqref{eq:M2}), the other zero;
\item use a first inert prime \(q_0\) with \(W(\gamma_0)\) a 3-cycle (\(\gamma_j\) rational Frobenius lift notation), giving an acyclic full local complex already by residual testing;
\item at each \(q_j,\ j>0\) require \(W(\gamma_j)=1\) and nonzero limiting Tate trace on all components and also on \(E_0\). Write \(\Phi_j=\lim\rho_P(\gamma_j),\ \alpha_j=\operatorname{tr}\Phi_j\) (also at \(j=0\), with unit trace). Then \(\det\Phi_j=-1\), \(\Phi_j^2-1=\alpha_j\Phi_j\), and \(\gamma_j^2\) is trivial on all ring class fields of \(K\) (inversion coset). One has a tame/Frobenius Koszul local complex on \(m_j,\alpha_j\Phi_j\), \(1+m_j=(1+u_j)^{b_j}\) for a unit \(b_j\in\mathbb Z_2^\times\), with terms
\[
                  P_R\longrightarrow (P_R)_s\oplus(P_R)_f
                      \longrightarrow P_R
\]
in degrees \(0,1,2\). Indeed use the two procyclic calculations of the diagram conventions (\(q_j^2\to1\)). For a summand line \(l\subset P_R\), impose the subcomplex \(U_j\) on \(l\to l_s\oplus(\Phi_j l)_f\to\Phi_j l\).
\end{itemize}

\begin{lemma}[The modified local conditions]\label{real:local-conditions}
For \(j>0\), the subcomplex \(U_j\) splits off and is exactly
self-orthogonal for conjugate local duality. The resulting complex
\(C\) has a perfect self-duality and a minimal square model
\(\delta_C:F\to G\) in degrees \(1,2\). The full strict problem also
has square amplitude \(1,2\).
\end{lemma}

\begin{proof}
The complex \(U_j\) and its degree-two-shifted dual, and the two analogues for a complementary line, are Koszul models of \(R/(\alpha_j,m_j)[-2]\) (regular sequence even over the order). Pairing morphisms are thus determined on \(H^2\), which can be checked on \(H^2,H^0\) pairing after reduction to that quotient (last-term cohomologies commute with reduction here). Then the local \(K\)-action is trivial, conjugate transport via \(\gamma_j\) is by \(\Phi_j\) up to scalar unit, and \(e(\Phi_j l,\Phi_j l)=0\) with the opposite entries for complementary lines units by the determinant form \(e\). Indeed the last term gives the tame-residue class up to unit trace normalization by the procyclic cochains and local duality at each quotient precision; Shapiro cup uses the invariant over \(K_v\), not twice that map. This proves nullhomotopy and quotient duality over \(R\).

Thus \(C\) has perfect self-duality and a minimal square model \(\delta_C:F\to G\) in degrees \(1,2\) by residual invariant vanishing over \(K\), injectivity of the local condition maps on \(H^0\), and duality. The full strict complex also has square amplitude \(1,2\) as before (dual invariant vanishing, and strict at \(w\) removes the global Euler contribution).
\end{proof}

\begin{proposition}[Residual transversality]\label{real:transversality}
The moving inert primes can be chosen so that
\[
                   (\det\delta_C)|_{t=\mathfrak m_A=0}\ne0             \tag{R1}\label{eq:R1}
\]
and likewise for the full strict determinants on \(P\) and \(T_2E_0\).
\end{proposition}

\begin{proof}
\emph{Frobenius weights at fixed places.}
At \(q_0\) require that the Frobenius exponents in \(u_0\) at \emph{every} fixed place over \(S\) are nonzero. Indeed take principal generators \(x_v\) for the \(h(K)\)-th ideal powers of one chosen place at each split prime. Along inert \(q_0\to-1\) the power symbol of order \(2^a\mid q_0+1\), \(a\to\infty\), kills rational residues and reads the translations on \(2^a\)-roots by \(\gamma_0^2\). Thus by ring class reciprocity and odd \(h(K)\) it detects the exponents via \(x_v\). Take lifts approaching \(c\eta,\ \eta\in G_K\) fixing the cyclotomic data and with cycle image giving \(W(c\eta)\) a cycle (compatible on the abelian intersection). On its square the translations are differences on conjugate numbers. Vary \(\eta\) within the deeper kernel of \(W\); the Kummer tuples for the \(x_v\) and their conjugates there with cyclotomic fixed have open image by independent valuations even after the finite extension (rational Kummer restriction survives the cyclotomic tower by a central nonidentity scalar). Thus all differences can be kept nonzero; apply Chebotarev. At the residual generic field in the \(u\)'s with \(t=0\) this gives vanishing local \(H^0\) at both dyadic places and acyclicity at fixed odd allowed support.

\emph{Parity of the residual problem.}
After each step including \(q_0\) the residual generic strict \(H^1\) there has even dimension \(h_{\rm res}\). Indeed compare with \(E_0\) over \(\mathbb Z_2[[\mathbf u]]\) at \(t=0\), using the fixed strict/full coordinates of \(C\) and full at moving support. Apply the parity test \eqref{eq:H1} first at the DVR of \((2)\), then along a characteristic-zero line germ at zero (start on an integral group-power line retaining generic ranks). At these localizations the moving inert terms are acyclic (inertia scalars or dummy-cycle at the first test, nonzero corresponding Tate traces at the characteristic-zero center). Thus the problem is self-dual there with square amplitude as above; pull back to the zero conditions at those places for \eqref{eq:H1}. At the characteristic-zero center use unramified inflation to remove the added places (unramified terms there acyclic too), getting the fixed-place problem. Here strict dimension is even by \(s_2(E_0/K)=1\) and the symmetric-plane dyadic parity flip above. This gives even residual generic dimension as claimed, since modified and full and the curve's and packet's problems agree there by acyclicity at varying and fixed odd allowed places.

\emph{Rank reduction.}
While \(h_{\rm res}\ge2\), add another \(q_j\) as specified with old residual rank-two evaluation in the finite plane of two independent old classes there. Initially omit its variable (unramified inflation). At \(u_j=0\) in the old residual generic field the new local action is trivial; take the pure singular plane instead of the finite one (with degree zero included, degree two omitted). These conditions are self-annihilating as in the square switch (pure cup calculation for \(q_j^2\to1\), conjugate transport by \(\gamma_j\) preserving the two axes). Thus rank two gives a drop by two. This change deforms in the residual coefficients to using the auxiliary local condition \([P_R\to(P_R)_s]\) in degrees \(0,1\), differential \(m_j\) (other Koszul differential zero there). Generically with \(u_j\) added it and the full and modified conditions give identical cohomology by acyclicity. Consequently the new \(h_{\rm res}\) is no greater than the switched dimension by the finite free models.

\emph{Simultaneous evaluation and nonzero traces.}
For existence of that choice use the product evaluations as in the imaginary case, taking \(\gamma_j\) approximating \(c\eta\) with \(\eta\) in the termwise product kernel \(\mathcal H\) of \(W\), cyclotomic and all old ring-field data over \(K\). Old residual strict classes inject into unrestricted crossed classes there on the product of the \(G_K\)'s. The conjugate induction (adding constant \(c\) to that product) is absolutely simple by absolute simplicity over \(K\) and distinct determinants of the scalar twists (infinite order already via \(u_0\)). Restriction to \(\mathcal H\) detects classes: the quotient contains a normal product of the \(C_3\)'s acting in \(W\) alone (other data abelian over \(K\)), whose diagonal \(C_3\) is central in that product without invariants, killing also its abstract \(H^1\). Thus joint evaluations on \(\mathcal H\) span full copies. Evaluation on \((c\eta)^2\) then projects via \(1+c\) to the original plane coordinates, still additively with full span. Some determinant is nonzero by the quadratic polarization test.

Perturb by products of squares within \(\mathcal H\) without changing the residual test. This can ensure all the nonzero limiting traces simultaneously. Indeed each characteristic-zero projection on \(G_K\) has Zariski closure containing \({\rm SL}_2\). To see it, irreducibility is already ensured by \(W\). If the identity group were solvable it would be toral (the normal unipotent radical acts trivially); potential scalars are excluded by the unequal Hodge--Tate weights. The two torus lines would then be preserved over a quadratic field for the full rational-base representation, necessarily \(L\) by residual irreducibility over other quadratic fields. But geometric one-dimensional characters over the real quadratic field have parallel weights: by local algebraicity of one-dimensional de Rham characters (classical abelian Hodge--Tate character theorem), their powers on small units at two, at the two embeddings combined, must be parallel by testing a power of a global infinite-order unit (inertia characters away from two have finite image). The two line characters interchanged over \(\mathbb Q\) would therefore have equal weights, impossible. This proves the Zariski assertion by the two-dimensional algebraic subgroup classification.

The constant second derived subgroup of \(G_K\) is available in \(\mathcal H\) termwise, since the excluded data have derived length at most two. Its projections and the groups generated by its squares still have closures Zariski containing \({\rm SL}_2\), by brackets/group commutators and squaring there. Thus on the simultaneous compact closure each trace-zero test on the coset of the limiting matrices has empty interior (every finite-index projection still Zariski contains the connected \({\rm SL}_2\)). Avoid them all and approximate by a product of squares. Chebotarev at increasing precision as in the evaluation lemma now achieves the requirements. Even parity and the successive drops prove \eqref{eq:R1}.
\end{proof}

Put \(s_C=\det\delta_C\). It is regular in \(R\), by
\eqref{eq:R1} and the embedding of the reduced coefficient order into
its integer components. Let \(K_b:G\to F^\vee\) be the degree-two
component of the compact pairing.

\Needspace{12\baselineskip}
\begin{lemma}[An integral alternating presentation]\label{real:pfaffian}
There is an integral matrix \(H\) such that
\begin{equation}\label{real:alternating-presentation}
 K'=K_b+\delta_C^tH\quad\hbox{is invertible},
 \qquad K'\delta_C\quad\hbox{is alternating}.
\end{equation}
The common rank of \(F,G\) is even. The element
\[
 \mathcal P=\operatorname{Pf}(K'\delta_C)
\]
is regular and remains nonzero after reduction in \eqref{eq:R1}.
Its Pfaffian adjugate, denoted \(\#(K'\delta_C)\), satisfies
\[
 \#(K'\delta_C)K'\delta_C=\mathcal P\,1.
\]
\end{lemma}

\begin{proof}
Set \(B=\delta_C^{-t}K_b\). Its reduction modulo integral matrices is
an alternating linking form with values in \(R[1/s_C]/R\).
For the diagonal assertion, test \(g\in G\) by
\(s_C\delta_C^{-1}g\bmod s_C\); its Bockstein from reduction modulo
\(s_C^2\) is represented by \(g\). After inverting all \(\alpha_j\),
\(j>0\), or all \(m_j\), \(j>0\), both the nonstandard conditions
and their full local complexes are acyclic. Replace them, and the
dummy-prime condition, by zero and apply \eqref{eq:H1}, with the
scalar clearing allowed there. The diagonal values of \(B\) are
therefore integral in both localizations. The products
\(\prod_{j>0}\alpha_j\) and \(\prod_{j>0}m_j\) form a regular
sequence, so these localizations intersect in \(R\). If there are no
such primes no intersection argument is needed. The same linking
pairing supplies the alternating off-diagonal relations.

Consequently \(B_{ii}\in R\) and \(B_{ij}+B_{ji}\in R\). Choose
\(H_{ii}=-B_{ii}\), and, for each unordered pair of distinct indices,
choose the two integral entries of \(H\) to cancel
\(B_{ij}+B_{ji}\). Then \(B+H\) is literally alternating. No division
by two is involved. Since
\[
 K'\delta_C=\delta_C^t(B+H)\delta_C,
\]
this proves alternation. Minimality gives \(\delta_C=0\) modulo the
maximal ideal, while perfectness makes \(K_b\) invertible there.
Thus \(K'\) is invertible. Since \(\det\delta_C\) is regular, the
alternating matrix has even size, and
\(\det(K'\delta_C)=\mathcal P^2\) proves the asserted regularity
and residual nonvanishing. The adjugate identity is the usual
integral polynomial identity for Pfaffians.
\end{proof}

\subsection{Cofactors and divisibility by all auxiliary traces}

Set \(t=0\). The arithmetic finite-model construction allows us to
use the matrices for \(P\) also on \(M=P\otimes_A D\), by tensoring
with \(D\). We also allow the flat integer coefficient extensions
used for the measures. Tensor factors are suppressed in the formulas.

\begin{lemma}[The cofactor lift]\label{real:cofactor-lift}
Let \(\eta\) be a closed degree-one input at \(w\), and let
\(\beta:C_w[-1]\to C\) be the strict-place boundary map. Put
\begin{equation}\label{real:lift-formula}
 y=\#(K'\delta_C)K'\beta\eta.
\end{equation}
There is a global cycle \(z\), obtained from \(y\) by integral maps
and an integral correction multiplied by \(\mathcal P\), such that
\[
 \delta_Cy=\mathcal P\,\beta\eta,
 \qquad \operatorname{loc}_w z=\mathcal P\eta
\]
in cohomology, after a consistent choice of sign. If a scalar \(c\)
clears \(\eta\), it also clears this construction of \(z\).
\end{lemma}

\begin{proof}
The Pfaffian adjugate identity implies
\[
 \delta_C\#(K'\delta_C)K'=\mathcal P\,1.
\]
It can first be checked after inverting \(s_C\), where it is the
inverse-matrix identity, and then holds integrally by regularity.
It remains true under every subsequent specialization, including one
where the specialized differential is singular.

Map back to the fiber model at \(w\), where \(\beta\) is the shifted
local inclusion, and include the comparison homotopies. Projection to
the problem relaxed at \(w\) gives a global cycle after subtracting
\(\mathcal P\) times the corresponding integral nullhomotopy value
on \(\eta\). The fiber equation gives the asserted localization.
Choose an unrestricted global minimal model with no degree-zero
term. All comparison maps and corrections are integral before being
applied to \(\eta\). Changing \(\beta\eta\) by \(\delta_Cv\) changes
\(y\) by \(\mathcal P v\), by the same adjugate identity, which also
accounts for changes of representatives. Linearity proves the
clearing assertion.
\end{proof}

We isolate the algebra that keeps the full product of the auxiliary
traces. In particular these traces need not be relatively prime.

\begin{lemma}[Pfaffian divisibility with repeated factors]\label{real:weighted-pfaffian}
Let \(R_0\) be a commutative ring, let \(\mathsf A,\mathsf B\) be
alternating matrices of even size, and let
\(\Delta=\operatorname{diag}(d_1,\ldots,d_{2r})\). Suppose
\[
 \mathsf A=\Delta\mathsf B\Delta,\qquad v=\Delta b,
 \qquad L=H\Delta
\]
with all matrices integral. The target of \(L\) may be any
\(R_0\)-module. If \(Q_\Delta=\prod_kd_k\), then
\begin{align}
 \operatorname{Pf}(\mathsf A)&=Q_\Delta\operatorname{Pf}(\mathsf B),
 \label{real:weighted-pfaffian-id}\\
 L\#\mathsf A\,v&=Q_\Delta H\#\mathsf B\,b.
 \label{real:weighted-adjugate-id}
\end{align}
Thus \(L\#\mathsf A\,v+\operatorname{Pf}(\mathsf A)c\) is divisible
by \(Q_\Delta\) for every integral \(c\). Neither regularity nor
coprimality of the \(d_k\) is required.
\end{lemma}

\begin{proof}
Every matching in the Pfaffian uses each index exactly once, proving
\eqref{real:weighted-pfaffian-id}. An entry of the Pfaffian adjugate
is a signed Pfaffian minor with two indices omitted. Restoring their
two diagonal factors gives, entrywise, the polynomial identity
\[
 \Delta\#\mathsf A\Delta=Q_\Delta\#\mathsf B.
\]
Multiplication by \(H\) and \(b\) proves
\eqref{real:weighted-adjugate-id}. These identities are polynomial,
so remain valid after every specialization and after tensoring with
any module.
\end{proof}

\begin{proposition}[Integral cofactor divisibility]\label{real:cofactor-divisibility}
Further specialize \(u_j=0\) for every \(j>0\), so that the
coefficient ring is \(R_1=A[[u_0]]\). For an integral input \(\eta\),
the global coordinates of the cycle in Lemma~\ref{real:cofactor-lift}
are divisible by
\[
 Q_{\rm aux}=\prod_{j>0}\alpha_j.\tag{R2}\label{eq:R2}
\]
For an input cleared by \(c\), the same assertion holds for \(cz\).
The assertion also holds with multiplicity coefficients \(D\).
\end{proposition}

\begin{proof}
\emph{A common local condition.}
At each specialized moving place, compare \(U_j\) with the finite
unramified low condition
\(P_{R_1}\to(P_{R_1})_f\). Their common subcondition is
\(l\to(\Phi_jl)_f\), in degrees \(0,1\). Quotienting by that common
condition leaves rank-one differential blocks \(\alpha_j\), in
degrees \(0,1\) for the low condition and \(1,2\) for \(U_j\).
Let \(C_B,C_0\) be the strict-pair problems with the low and common
conditions. The low problem is exactly self-dual: the unramified cup
is null, and the residual local test gives quotient duality. It has
a minimal square model \(\delta_B:F_B\to G_B\) in degrees \(1,2\).
Its size is even. To see this at the residual closed point, remove
the specialized places by unramified inflation and remove the
acyclic dummy place. The resulting residual complex agrees with the
elliptic fixed-support problem over \(\mathbb Z_2\). Identity
\eqref{eq:H1} and the characteristic-zero strict-pair parity used in
Proposition~\ref{real:transversality} give even closed-point dimension.

Take the fiber of the map from \(C_B\) to the low quotients, and then
extend by the high quotients contributed by the \(U_j\). This gives
an integral model for the specialized \(C\):
\begin{equation}\label{real:block-model}
 \begin{gathered}
 \delta_L:F_B\oplus R_1^m\oplus R_1^m
       \longrightarrow G_B\oplus R_1^m\oplus R_1^m,\\
 \delta_L=
 \begin{pmatrix}
  \delta_B&0&a\\
  e_0&T_\alpha&k_0\\
  0&0&T_\alpha
 \end{pmatrix},\qquad
 T_\alpha=\operatorname{diag}(\alpha_j)_{j>0}.
 \end{gathered}
\end{equation}
Here \(m\) is the number of places and signs have been absorbed in
the bases. The first two blocks model \(C_0\) and project to \(C_B\).
All maps in this construction are the local-condition comparisons.

\emph{Transport of the alternating presentation.}
The model \eqref{real:block-model} splits over the local ring into a
minimal part and unit differential disks. The derived comparison with
the original specialized model, which is still minimal, is a
degreewise isomorphism on the minimal parts: a homotopy equivalence
of minimal free complexes is an isomorphism after reduction and hence
before reduction. Pull back \(K'\) there. Both minimal sizes are even,
so the remaining number of unit ranks is even; equip those disks
with a unimodular hyperbolic alternating pairing. We obtain an
invertible \(K_L\) such that \(K_L\delta_L\) is alternating.
Moreover \(K_L\) differs from the actual compact-cup component by
\(\delta_L^t\) times an integral matrix. This is the form of the
linking modification and of a homotopy difference on two-term
models; the unit disks are contractible.

Pfaffian and adjugate congruence formulas transport the lift, up to a
scalar unit and \(\mathcal P\)-multiples of integral corrections.
The disk block has unit Pfaffian. Projection to the minimal part
therefore transports the lift formula itself, without inverting the
specialized differential. Boundary entries that differ by a
differential change the lift by a \(\mathcal P\)-multiple, as in
Lemma~\ref{real:cofactor-lift}. The transferred global homotopies have
the same property, because \(\delta_Cy=\mathcal P\beta\eta\) and the
global model has no degree-zero term. By naturality the boundary
entry in the large model is the image of the common boundary in
\(C_0^2\).

\emph{The factors on rows, inputs, and projections.}
We verify the three factorizations needed for
Lemma~\ref{real:weighted-pfaffian}: one for the alternating matrix,
one for the boundary input, and one for projection to global cochains.
Let \(r_i\) be a low extra column and \(e_i\) its low row. Then
\(\delta_Lr_i=\alpha_i e_i\). For every \(g\in C_0^2\), the
\(r_i\)-coordinate of \(K_Lg\) is divisible by \(\alpha_i\).
Indeed the compact cup restricted to \(C_0\), whether pulled back
through this model or through \(C_B\), agrees up to homotopy. Keep
the unchanged isotropy nullhomotopies identical. On each changed
common condition, which has free terms only in degrees \(0,1\),
their difference has no degree-three tensor component and hence no
ambiguity. The projection to \(C_B\) kills \(r_i\), and every
homotopy or linking-modification term on \(r_i\otimes g\) factors
through \(\delta_Lr_i=\alpha_i e_i\).

Set \(\mathsf A=K_L\delta_L\). All entries incident to \(r_i\)
therefore contain \(\alpha_i\). More precisely, for two distinct
low indices,
\[
 \mathsf A_{r_i,r_j}
   =\alpha_j(K_Le_j)_{r_i}
   \in\alpha_i\alpha_jR_1.
\]
This product assertion is essential when traces share factors. Let
\(\Delta\) have entry \(\alpha_i\) at \(r_i\) and entry one at
all other indices. The preceding factorizations give an alternating
integral \(\mathsf B\) with \(\mathsf A=\Delta\mathsf B\Delta\).
The boundary vector \(v=K_L\beta\eta\) is \(\Delta b\), since
\(\beta\eta\in C_0^2\). Finally the underlying global projection
has the form \(H\Delta\): on \(C_0\) it factors through \(C_B\)
up to homotopy, and, since the target has no degree-zero term, its
value on \(r_i\) is an integral map applied to
\(\delta_Lr_i=\alpha_i e_i\).

Lemma~\ref{real:weighted-pfaffian} now shows that the projected
adjugate lift and every Pfaffian correction contain the full product
\(Q_{\rm aux}\). This proves \eqref{eq:R2}. All maps and
factorizations are integral before tensoring with \(D\); they
therefore remain factorizations afterward. The constant traces act
injectively on that multiplicity lattice by rational semisimplicity
and integer torsion-freeness. Applying the argument to \(c\eta\)
proves the clearing assertion.
\end{proof}

\begin{lemma}[Inflation at the center]\label{real:center-inflation}
At the all-zero center, inflation from fixed support identifies the
integral degree-one global groups, including with \(M\) coefficients,
and is compatible with the transport actions.
\end{lemma}

\begin{proof}
At every added prime the singular block in unramified inflation is
an injective Frobenius differential in degrees \(1,2\). It is the
limiting trace times an invertible matrix, since \(q_j^2\to1\).
The same injectivity holds with multiplicity coefficients.
The comparison triangles therefore identify degree-one cohomology
integrally; the finite diagrams preserve the actions.
\end{proof}

\subsection{Paired and Pfaffian divisibilities}

The passage from primitive components back to the Hecke order will
use the following elementary form of Weierstrass division.

\begin{lemma}[Divisibility over a coefficient order]\label{real:order-descent}
Let \(B\hookrightarrow\prod_\lambda B_\lambda\) be an injective map
of complete Noetherian local rings, with each component map local.
Put \(R_B=B[[\mathbf z,x]]\) and
\(R_\lambda=B_\lambda[[\mathbf z,x]]\). Suppose \(f\in R_B\) is
\(x\)-distinguished up to a unit, of degree \(d\). Then
\begin{equation}\label{real:quotient-injection}
 R_B/(f)\longrightarrow\prod_\lambda R_\lambda/(f_\lambda)
\end{equation}
is injective. If \(B\) and all \(B_\lambda\) are two-torsion-free,
divisibility of an integral series by \(f_\lambda\) in each
\(R_\lambda[1/2]\) implies divisibility already in \(R_B\).

If initially \(f\in B[[t,\mathbf u]]\) and
\(f|_{t=\mathfrak m_B=0}\ne0\), the distinguished hypothesis can be
achieved by a change of tame variables alone, leaving \(t\) fixed.
\end{lemma}

\begin{proof}
Write \(f=uF\), where \(u\) is a unit and \(F\) is monic and
distinguished of degree \(d\). For \(g\in R_B\), divide it by \(F\): write
\(g=Fq_0+r\) with \(\deg_xr<d\), and set \(q=u^{-1}q_0\).
Then \(g=fq+r\). Since the component maps are local,
\(F_\lambda\) remains distinguished of the same degree and
\(u_\lambda\) remains a unit. If \(g_\lambda\in(f_\lambda)\),
uniqueness of division gives \(r_\lambda=0\). Injectivity on the
coefficient rings gives \(r=0\), proving
\eqref{real:quotient-injection}. Moreover
\[
 R_B/(f)\simeq B[[\mathbf z]]^{\oplus d}
\]
as a module, and the same holds on every component. These quotients
are two-torsion-free under the stated hypothesis, so a class that
vanishes after inverting two already vanishes integrally.

For the final assertion, let \(h\) be the nonzero residual series in
the tame variables. Choose the lexicographically least exponent
\((a_1,\ldots,a_s)\) in its support. Take \(N_s=1\) and, successively,
\(N_i>\sum_{j>i}a_jN_j\). The chosen exponent is the unique one of
least weighted degree for these weights. Thus
\(h(x^{N_1},\ldots,x^{N_{s-1}},x)\ne0\). The triangular change
\(u_i\mapsto u_i+x^{N_i}\) for \(i<s\), with \(u_s=x\), makes
\(f\) distinguished up to a unit over
\(B[[t,u_1,\ldots,u_{s-1}]]\). This also works over the finite
residue field \(\mathbb F_2\).
\end{proof}

Take determinants \(s_F,s_*\) of the full strict problems for \(P,T_2E_0\) respectively. Use the paired measures \eqref{eq:M1}--\eqref{eq:M2} over \(\mathcal V^*\) in the same limits, odd conductor the product of the current inert primes including the dummy, with canonical connected level at 2 and depleting exactly at \(S\). Take base centers by simultaneous descending cyclic quotients from a conductor-one point with transported level; use common orientations and underlying points/parameters for comparisons. Write \(b^\pm\) for the two opposite odd orientations, and use a subscript \(*\) for the elliptic \(E_0\) normalized form. For the packet use the \(A\)-tuple of primitive normalized forms (of their respective true levels). Indeed their fully depleted antiderivative on common raised level has expansion
\[
                         \sum_{(m,\prod_{q\in S}q)=1} T_m q_{\rm Tate}^m/m,
\]
so gives integral \(A\)-coefficients by the expansion comparison at \eqref{eq:M1} with a basis lattice of \(A\) (first allow bounded coefficient denominator). Likewise the measures agree residually with the elliptic ones. Extend power-series coefficients by the unramified integer limit as before; the extended coefficient order for \(A\) is still local with the corresponding residual map (\(A\) itself has residue \(\mathbb F_2\)).

\begin{proposition}[Paired and Pfaffian division]\label{real:division}
The quotients
\[
 V_A=b^+ b^- /s_F,\qquad V_*=b_*^+ b_*^-/s_*,\qquad R_b=b^+/\mathcal P
                                                                        \tag{R3}\label{eq:R3}
\]
are integral in the power-series rings just described. If \(V_*\)
is a nonunit, then \(R_b\) is a nonunit.
\end{proposition}

\begin{proof}
\emph{The full paired quotients.}
First take a high finite character \(\theta\) in \(t\) and use either \(E_0\) or a characteristic-zero primitive component \(f\) of \(A\). The divisibility after inverting 2 for the full quotients is exactly the \eqref{eq:B1}--\eqref{eq:D2} comparison, at all height-one tests of the multi-tame-variable arithmetic ring with 2 inverted. We check the application here.
\begin{itemize}

\item Strict generic acyclicity is by \eqref{eq:R1}. The added inert locals are contractible at these tests by the nonzero limiting traces on each problem (and the two procyclic blocks); at odd fixed \(q\) the unramified replacement accounts for \(P_q^f(\xi_q)P_q^f(\xi_q^{-1})\) by the singular-block calculation before \eqref{eq:B1}, with \(P_q^f\) the primitive Euler polynomial in the normalization at \eqref{eq:M1}. Both are generically nonzero by the nonconstant powers in \(u_0\).
\item At 2 use the several-variable high-character Shapiro/Kummer calculation of the local comparisons: Frobenius has a nonzero \(u_0\)-exponent, with all the \(u\)-actions unramified there, giving the indicated flat substitution. Base ring fields locally contain each required fixed unramified layer eventually and the \(\theta\)-field; ramification stays bounded at this fixed character depth. Thus log identification and transfer use the unnormalized weighted sums as at \eqref{eq:B1}; use full modular abelian factors then project, with fixed coefficient denominators.
\item Use extra derivative primes and the points at the primitive level exactly as in \eqref{eq:K} and \eqref{eq:D2}, with base conductor \(2^a\prod q_j\) here for fixed dyadic depth \(a\). The relative inertia description holds by the order class sequence with our unit hypothesis; base ramification at fixed places stays bounded (composita as in \eqref{eq:K}), and no Kummer restriction at the existing moving inert primes is required. Absolute simplicity of the Tate plane over \(K\) already follows from \(W\), and \(\chi^4\ne1\) by high dyadic inertia. Thus the full-factor homothety test, rank-two addresses, and comparisons via \eqref{eq:K} prove the nonnegative valuation in \eqref{eq:D2} at each height-one DVR if the measure product is nonzero. Primitive orientation sums differ by Atkin--Lehner as in \eqref{eq:M2}, giving \eqref{eq:B1} with the two Euler translations just stated. Consequently over the arithmetic ring the square of the localization coefficient is divisible by the odd-unramified-replaced strict determinant, by all those valuations; after extended log identification this gives the full divisibility for \(b_f^+ b_f^-\) by \((s_F)_f\), and analogously on \(E_0\).
\end{itemize}

For each sufficiently high \(\theta\), residual regularity in
\eqref{eq:R1} permits distinguished division in a tame coordinate.
The quotient by that distinguished denominator is two-torsion-free,
so the horizontal divisibility gives integral divisibility at the
character. Apply Weierstrass division before evaluating \(\theta\).
Its bounded remainder vanishes at all sufficiently high characters,
and hence vanishes identically by the bounded-series character test.
Lemma~\ref{real:order-descent} checks this conclusion on the integer
components and descends it to the coefficient order. Therefore
\(V_A,V_*\) are integral.

\emph{Comparison with the Pfaffian.}
For the \(A\) problem, work before specialization on a primitive
component \(f\), and put \(S^\circ=S\setminus\{2\}\). There are
\emph{integral power-series units} \(\varepsilon_f,\eta_f\) such that
\begin{align}
 b_f^-\prod_{q\in S^\circ}P_q^f(\xi_q)
 &=\varepsilon_f b_f^+
       \prod_{q\in S^\circ}P_q^f(\xi_q^{-1}),
 \label{real:opposite-orientations}\\
 \frac{(s_F)_f}{\mathcal P_f^2}
 &=\eta_f\prod_{q\in S^\circ}
       \frac{P_q^f(\xi_q^{-1})}{P_q^f(\xi_q)}.
 \label{real:determinant-ratio}
\end{align}
Indeed in \eqref{eq:M2} the primitive Atkin--Lehner at the full odd level part relates the sums by sign and translation by the relevant oriented prime ideals of that part (group powers in the variables, preserving canonical level at 2). The Gauss and parameter bases use the same underlying points and choices. Thus the cross identity with this unit holds at all sufficiently high characters, which suffices by boundedness.

For the determinant the moving modified-vs-full complementary rank-one Koszul blocks have determinant units (regular sequence \(m_j,\alpha_j\), acyclic in codimension one on integer components). At each fixed odd pair one adds the full complex at the place of \(\xi_q^{-1}\). At level-good \(q\) the singular over finite Frobenius determinant gives exactly the ratio up to unit. At minimal ramified transposition inertia, conductor one means Steinberg and the inertia difference has a unit pivot by the residual action, giving free invariant and coinvariant lines integrally with Frobenius \(q a_q,a_q\), \(a_q=\pm1\); thus again the stated ratio. At inertia containing \(C_3\) there is acyclicity by residual testing and no Euler line. And at the level-raising prime \(\ell_0\) the full complex is acyclic integrally by residual Frobenius, and the Euler polynomials units (level-good with odd trace, or exponent two with no invariant line; a line would force unipotent inertia of conductor at most one). These exhaust the cases by the level construction. Pfaffian-square differs by unit from the modified determinant, proving the identity by triangles.

Combining these identities gives
\[
 V_{A,f}=\varepsilon_f\eta_f^{-1}
       \left(\frac{b_f^+}{\mathcal P_f}\right)^2.
\]
Normality on each integer component makes \(b_f^+/\mathcal P_f\)
integral. Residual regularity of \(\mathcal P\), established in
Lemma~\ref{real:pfaffian}, and Lemma~\ref{real:order-descent} then
give \(R_b\) in the coefficient order itself. Merely knowing
integrality on its normalization would not suffice without this
quotient-injectivity argument.

\emph{Nonunits.}
The numerator congruences and identical residual strict diagrams
show that \(V_A\) and \(V_*\) agree modulo the residual map up to a
unit: cancel their common nonzero residual denominator. Thus a
nonunit \(V_*\) makes \(V_A\) a nonunit. If \(R_b\) were a unit,
the displayed component identities, whose multipliers are integral
units, would make \(V_A\) a unit on every integer component. Every
maximal ideal of those integral components contracts to the unique
maximal ideal of the coefficient order; an element in that maximal
ideal cannot become a unit on a component. This contradiction proves
the nonunit assertion and completes \eqref{eq:R3}.
\end{proof}



It remains to construct an integral \(\sigma\)-fixed global class
\(p\) with \(Y=R_b(0)p\). The next two comparisons have different
roles: the horizontal logarithm identifies the global lift after
inverting two, while the integral clearing calculation controls its
remaining dyadic denominator.

\subsection{Horizontal logarithms at the anticyclotomic center}

In what follows set \(t=0\) first, keeping all the \(\mathbf u\). Put
\[
                   x=\chi(\mathrm{Fr}_w),\qquad \xi_2=x^{-1}.
\]
Both dyadic characters are now unramified, with opposite exponents; the \(u_0\)-exponent of \(x\) is nonzero. Work for now after 2-inversion on the plane \(T\) of a primitive \(f\) occurring in \(P\), realized also by its modular abelian variety \(\mathcal A_f\). True level at two is good or of exponent one as above. Fixed isogeny/projector and lattice denominators are allowed in this paragraph and the following \emph{horizontal} calculations. Use the Lie line with log coordinate pulling back (by the chosen projection from the primitive level) to \(f\,dq_{\rm Tate}/q_{\rm Tate}\). We work over DVRs of height one on the 2-inverted integer power-series ring after extending constants as on the measure side. Finite further coefficient extensions to split the projections/polynomials can be used, component by component (orders then scale by ramification at a test). Also allow the DVR of a characteristic-zero group-power line germ at zero retaining nonconstant \(x\).

\begin{lemma}[The horizontal local logarithm]\label{real:local-log}
There are perfect local maps \(U_w,U_{\bar w}\), exact orthogonals
over the test DVR. At the generic field they are lines in degree one injecting in full local cohomology (of rank two in degree one only). At \(w\) we identify that generic line by weighted log and claim
\[
                     d(U_w,e_{\log})=v(P_2^f(x^{-1}))                 \tag{R4}\label{eq:R4}
\]
for the generic basis of log coordinate one. On the test residue field the maps inject in degree zero with no negative degrees in the conditions. At the characteristic-zero center on the line germ they give the
ordinary finite Kummer lines; full local cohomology there is also
concentrated in degree one.
\end{lemma}

\begin{proof}
\emph{The unramified parameter and the coefficient plane.}
As in the high-character calculation, do unramified Shapiro first on \(D_j/\mathbb Q_2\) of degree \(2^j\), with universal Frobenius value \(1+s\); then use \(1+s=x\), or the opposite for the conjugate condition. This agrees with the local finite models by stabilization at every precision. The log identification uses \(\sum_k(1+s)^k F^k\) on log points, in cyclic-quotient notation with \(F\) arithmetic Frobenius, evaluated at our fixed embedding before extending on coefficients. Frobenius on the underlying point system itself thus translates the sums by \((1+s)^{-1}\).

At good level \(\mathcal A_f\) has good reduction by modular local compatibility. The relation of special Frobenius with Hecke at 2 gives roots of \(Z^2-a_2 Z+2\) on the chosen coefficient projection. We use the standard Newton and connected/étale descriptions of the good-reduction 2-divisible group: unit-root rank is étale rank (Frobenius invertible on that part), with arithmetic action on the lifted étale Tate, and the two slopes here sum to one. Indeed use the Frobenius polynomial relation on reduction by good Eichler--Shimura, on the plane in Dieudonné theory; polarization with the good Hecke field self-adjoint also gives slope symmetry on that plane. In the case of a unit root \(\lambda\) this gives
\[
     0\to T^+\to T\to T^-\to0,\qquad
              T^-=\eta_\lambda,\qquad T^+=\eta_\lambda^{-1}(1)
\]
by taking the étale quotient and using the determinant. Here \(\eta_\lambda\) is the unramified character of arithmetic value \(\lambda\), which in this good case isn't a root of unity by purity. Without a unit root the étale projection is zero and both slopes are strictly between zero and one.

At true exponent one the primitive factor has semistable purely toric reduction at 2, by semistable reduction and the toric rank criterion on auxiliary Tate modules (Steinberg compatibility, unipotent inertia with monodromy rank half). Use Tate--Raynaud uniformization. We get the same sequence, with \(\lambda=a_2=\pm1\): the periods and torus cocharacters are identified rationally by valuation, their Frobenius sign on the factor being the Steinberg quotient sign, as also computed on the inertia coinvariants of auxiliary Tate. The torus splits over a fixed unramified extension (the sign over at most the quadratic extension). After unramified splitting, valuation of the period extension on the plane is nonzero, by that full-rank valuation isomorphism. Throughout, pair the two graded pieces by the determinant form (polarization plane form up to scalar constant).

In the cases with \(T^\pm\), use provisionally
\[
  U_w^\circ=\mathrm{fib}\big(C_w(T^+\chi)\to
                               C_{\rm ur}(T^-\chi^{-1})^\vee[-2]\big)
                   \ \longrightarrow\ C_w(T\chi),
\]
using local duality and the ordinary two-term unramified complex in degrees \(0,1\). This condition itself (as a complex) is a free log line in degree one over the horizontal test ring. Indeed unramified inflation on an unramified coefficient includes \(H^0\) isomorphically and injects on \(H^1\) on every fiber; the cone has no degrees below one already by the residual lattice test over the universal ring. Also \(H^0\) on \(T^+\chi\) vanishes on every characteristic-zero field test by cyclotomic inertia (degree-zero evaluations inject after base change by the finite cochain models). Hence duality and local Euler characteristic give just a line in degree one on these fields for \(U_w^\circ\). Before field tests, \(H^1(U_w^\circ)=H^1(C_w(T^+\chi))\) since the dual target has injective Frobenius differential in degrees \(1,2\). These checks also apply over the 2-inverted universal unramified ring before substitution.

\emph{Uniform logarithm bounds.}
Weighted logarithm identifies this line up to a unit of the ring with
two inverted and constants extended. This is a horizontal statement;
it makes no assertion that a fixed two-power lattice denominator is
an integral unit. The required uniform bounds are as follows. In good reduction let \(T^\circ\) denote the formal/connected Tate of the actual variety. Formal Kummer embeds the completed formal points at each \(D_j\) in \(H^1(D_j,T^\circ)\) with torsion-free cokernel (the cokernel injects in a Tate module by Kummer exactness). Indeed formal points over the algebraic closure are 2-divisible: divisions in the variety can be adjusted to reduction zero using geometric torsion, surjecting onto special-fiber torsion points by finite flatness. The ordinary \(T^+\)-projection of that inclusion spans rationally at every \(D_j\): formal Lie has rank one on the embedding, and \(H^1\) on \(T^+\) has the corresponding dimension by local Euler characteristic and \(\lambda^{2^j}\ne1\). Thus compatible cohomology tuples on a lattice of \(T^+\) come from formal Kummer up to a \textbf{fixed} projection/lattice denominator. More explicitly, map back after such a fixed multiple into \(T^\circ\) with scalar extension; the image lands rationally in formal Kummer so integrally there by saturation before projecting, and preimages there are unique. Formal logs are bounded above with torsion kernel of bounded exponent over all the unramified layers, and cover a uniform small additive Lie lattice, by the formal log and exponential on a small ball (a bounded power of 2 suffices to enter it).

In the toric case compute instead with the torus Kummers for \(T^+\). In the split layers multiplicative Kummer gives all the degree-one groups; the valuation coordinates die in the inverse norm limit (multiplication by the growing degrees), and the units give the same formal estimates. On our projection the cocharacter action is just the sign as above; fixed unramified splitting or lattice descent thus costs only constants. The Lie map of uniformization is an isomorphism; a fixed multiple of formal abelian points also lifts to small torus points, compatibly with trace by their log coordinates.

Pass by compact inverse cohomology at integral arithmetic precision. Trace-compatible integral normal generators \(z_j\in\mathcal O_{D_j}\) with unit total trace give free inverse additive lattices over the arithmetic unramified power-series ring, as in the high-character comparison. Thus after fixed denominators, \(H^1\) of our pointed condition over that ring with 2 inverted is identified by log systems with the inverse Lie line; the bounds above and compact exact limits control torsion and indices. The resolvent of the \(z_j\), namely the inverse group sum with weights \((1+s)^k\), is an integer-series unit over extended constants by its augmentation. This gives the asserted unit log evaluation (also after substitution; the free line complex pulls back). In particular a compatible formal Kummer system from \(\mathcal A_f\) belongs to the condition generically, with log evaluated this way.

\emph{The horizontal divisor and specialization.}
Except in the good ordinary case at \(x=\lambda^{-1}\) on the test residue field, take \(U_w=U_w^\circ\), and take \(U_{\bar w}\) its exact orthogonal by the pairing fiber. In that exceptional case take \(U_{\bar w}=U_{\bar w}^\circ\) analogously instead, and \(U_w\) its exact orthogonal. These choices have the isotropy nullhomotopies by construction. The orthogonal fits in a triangle from the full \(C(T^+\otimes\text{scalar})\) at that place with quotient \(C_{\rm ur}(T^-\otimes\text{scalar})\): it is the pullback of that unramified map along the projection by duality of the filtration. Hence there is degree-zero injectivity on all fibers. In the exceptional case \(C_w(T^+\chi)\simeq U_w^\circ\) at this DVR since \(\lambda\ne\lambda^{-1}\); the quotient contributes \(+v(1-\lambda x)\) to \eqref{eq:R4}. At all other tests \eqref{eq:R4} needs exponent zero (toric polynomial \(1-a_2/(2x)\); in good ordinary case use \(2x^2 P_2^f(x^{-1})=(1-\lambda x)(1-2x/\lambda)\)). At the characteristic-zero center the good ordinary assertion follows by the formal Kummer description. For toric \(\lambda=1\) at the center, \(U_w^\circ\) uses the units line via removing the valuation by the dual unramified cup. It injects as usual finite Kummer despite the period relation by the nonzero period valuation; there are no Tate invariants. For \(\lambda=-1\) there is no exceptional phenomenon there and the torus line likewise gives the Kummer line. The orthogonal then specializes to opposite finite by abelian local duality.

\emph{The nonordinary case.}
In good nonordinary reduction use the inverse Kummer line at \(w\) as in the earlier projected log comparisons, then its exact orthogonal. Here there is no étale part rationally on that projection; costs from reduction are bounded uniformly by a projector denominator. Indeed the geometric étale 2-power points on reduction are divisible, so an integral multiple of the projector killing their rational Tate module annihilates the whole geometric 2-primary point group (after the coefficient extension); use endomorphisms before and after reduction. Thus after this bounded loss over any \(D_j\) we can use formal points for upper log and torsion estimates, and small-log lower bounds are uniform as well. Integral Kummer and dual on the full variety as in the high-character proof then give, after 2-inversion and projection, two ends of the \(H^1\) sequence free over the unramified-variable ring (additive normal bases by log and trace dual log), with \(H^2\) zero by the torsion bounds. Thus the local complexes are themselves just split degree-one planes with direct Kummer lines, and evaluation is a log isomorphism on the projected inverse Kummer line. It specializes to Kummer at zero by the small-log generators, so its orthogonal likewise does. Both polynomial roots (in the Frobenius eigenvalue convention) have positive slope, giving exponent zero required in \eqref{eq:R4}. This completes the local comparison.
\end{proof}

\subsection{An input with an integral clearing factor}

At \(t=0\), still, let \(Y(\mathbf u)\) be the Kummer of the positive-orientation conductor-\(\prod q_j\) trace on \(J\) with the \(\chi\)-weights (stage/limit notation). Use the level \(n\), canonical connected choice at 2, and subtract the \(\infty\)-cusp on the original point before weighted corestriction. Project to \(M\). Use the finite models of the cofactor calculation with constants extended as on the measure side. In the corresponding projected Lie of \(J\) set
\[
 \mathcal L=\prod_{q\in S} B_q^{-1} a_1^A,\qquad
 B_q=1-(\mathcal T_q/q)\xi_q+(\epsilon_{q,n}/q)\xi_q^2 ,
\]
where \(\mathcal T_q=T_q\) if \(q\nmid n\), \(U_q\) otherwise, \(\epsilon_{q,n}=1,0\) respectively. These are Lie operators with the usual dual action on differentials. \(a_1\) is the tangent vector extracting the first expansion coefficient of a pulled-back differential (coefficient of \(dq_{\rm Tate}\)), with superscript denoting Hecke projection.

\begin{lemma}[An integrally controlled logarithmic input]\label{real:input-clearing}
There is a closed input \(\zeta\) in the degree-one \(w\)-model,
in the generic Kummer lines of Lemma~\ref{real:local-log} on all
primitive projections, whose weighted logarithm is \(\mathcal L\).
There is a scalar clearing series \(c\) for \(\zeta\) such that
\[
 \overline c(u_0)
   :=c\big|_{2=u_j=0\ (j>0)}\ne0.
\]
The construction and clearing assertion are compatible with the
multiplicity comparison \(M=P\otimes_A D\).
\end{lemma}

\begin{proof}
\emph{The N\'eron tangent lattice.}
The tangent vector \(a_1\) belongs to the N\'eron Lie lattice. To see
this, take a smooth cusp chart above \(\infty\) with fine full
auxiliary prime-to-two tame level. The only possible extra level at
two is the canonical multiplicative subgroup of the Tate curve,
since \(2\|n\) when two divides \(n\). The cusp model for
generalized elliptic curves with prime-to-two and canonical
\(\Gamma_0(2)\) level is smooth there, over unramified constants,
with parameter width prime to two
\cite{KatzMazur85,DeligneRapoport73}. The cusp Abel map extends from
the generic fiber by the N\'eron mapping property. Pulling back
invariant differentials proves the assertion, since the width factor
is a unit. The Hecke operators preserve the N\'eron lattice, and the
anemic idempotent projection is integral.

On the projected lattice, the scalar series
\(\det(2B_2)\) and \(\det B_q\) for odd \(q\in S\) remain nonzero
on the residual \(u_0\)-line. Their arguments are nonconstant by the
choice of Frobenius weights. At two, use the monic quadratic term if
\(2\nmid n\); if \(2\mid n\), use \(\mathcal T_2^2=1\), by the
Steinberg assertion of Lemma~\ref{real:packet}. Thus
\[
 c_0=\det(2B_2)\prod_{q\in S\setminus\{2\}}\det B_q
\]
is a scalar series with nonzero residual restriction, and the
adjugate formulas show
\(c_0\mathcal L\in2\operatorname{Lie}_J\). The factor two is
explicit: \(B_2^{-1}=2(2B_2)^{-1}\).

\emph{Formal logarithms with a residual nonzero multiplier.}
Choose smooth N\'eron formal coordinates over \(\mathbb Z_2\) and
an integral tangent basis. For unramified integer coordinates \(y\),
\begin{equation}\label{real:formal-log-mod-two}
 \frac{\log(2y)}2\equiv y+H F(y)\pmod2,
\end{equation}
where \(H\) is a constant residual matrix and \(F\) is arithmetic
Frobenius. Here \(2y\) means substitution of the doubled coordinate
vector in the formal logarithm. Indeed its derivatives are integral.
For a monomial of total degree \(k\), the denominator valuation is
bounded by the valuation of the greatest common divisor of its
exponents; after the substitution and division by two, its valuation
is at least \(k-1-v_2(k)\). Only the linear and pure quadratic terms
can survive modulo two. The pure squares give \(F(y)\) on
unramified residue fields, proving \eqref{real:formal-log-mod-two}.

Let \(h(Z)\in\mathbb Z_2[Z]\) lift \(\det(1+HZ)\). Apply the
adjugate of \(1+HF\) modulo two, and then the logarithm isomorphism
on the small formal ball with coordinates in \(4\mathcal O_{D_j}\).
It follows that the logarithms of formal points cover
\[
 2h(F)\operatorname{Lie}_J(\mathcal O_{D_j})
\]
at every unramified layer \(D_j\). For a trace-compatible integral
normal generator \(z_j\) and an integral tangent vector \(a\),
choose preimages of \(2h(F)z_ja\). These preimages can be chosen
norm-compatibly: any finite set of norm constraints is solved by
choosing at the highest layer and norming down, and the fibers are
compact. Compactness then solves all constraints simultaneously.

Kummer theory and unramified Shapiro give cycles with logarithmic
coordinates \(2h(x^{-1})a\), multiplied by the unit normal
resolvent of Lemma~\ref{real:local-log}. This also holds on finite
character tests. Since \(\overline h\) has constant term one and
\(x\) is nonconstant on the residual \(u_0\)-line,
\(\overline h(x^{-1})\) is not the zero series. Combine the cycles
linearly with \(c_0\mathcal L\in2\operatorname{Lie}_J\), and divide
by the scalar product \(c_0h(x^{-1})\) and by the unit resolvent.
This constructs \(\zeta\) with a clearing series having nonzero
residual restriction. The unrestricted Tate cochains, their maps to
the finite models, and the tensor comparison with \(D\) are used
before these scalar combinations, so the assertion also holds on
\(M\).
\end{proof}

\subsection{The global class and its integral specialization}

Lift \(\zeta\) by Lemma~\ref{real:cofactor-lift} to a global class
\(z\), with the sign chosen so that its localization is
\(\mathcal P\zeta\). The defining logarithm of \(\zeta\) will
identify this lift with the weighted Heegner class as follows.

\begin{proposition}[Identification with the Heegner class]\label{real:heegner-identity}
Over the generic fields,
\[
                             Y(\mathbf u)=R_b z .                     \tag{R5}\label{eq:R5}
\]
\end{proposition}

\begin{proof}
We first verify the logarithmic identity at \(w\).

For each \(f\) at true level \(N_f\mid n\) use the projected primitive-quotient pushforwards \(\pi_i\) from degeneracy maps, just \(i=0\) unless \(n=N_f\ell_0^2,\ \ell_0\nmid N_f\), and \(0\le i\le2\) then (argument multiplied by \(\ell_0^i\) in analytic uniformization). They give rational coordinates on \(M\) by newform theory. Normalize with the same primitive log for each, i.e. pullback on differentials
\[
                  \omega_i=\ell_0^i f(q_{\rm Tate}^{\ell_0^i})
                                          dq_{\rm Tate}/q_{\rm Tate}.
\]
Then exactly
\[
                 \pi_i\mathcal L=\xi_{\ell_0}^i\Big/
                                         \prod_{q\in S}P_q^f(\xi_q)
\]
in these coordinates (numerator one for \(i=0\)). At unchanged good or minimal levels use the eigenlaw. In the old case evaluate \((1-U_{\ell_0}\xi_{\ell_0}/\ell_0)^{-1}\) on differentials with a formal symbol first, using \(a_1 U_{\ell_0}^k(\omega_i)=a_{\ell_0^k}(\omega_i)\); good Hecke recursion gives exactly the displayed rational multiplier.

On the conductor centers \(\pi_0\) uses the CM level at \(N_f\); the other projections use \(i\) successive oriented \(\ell_0\)-translations before the primitive pushforward, multiplying log weights by \(\xi_{\ell_0}^i\), with the same orientations as in \eqref{eq:M2}. Also the connected \(V_2\)-quotient there is arithmetic \(F=\mathrm{Fr}_w\), by Frobenius on the ordinary reductions including the transported tame data and uniqueness of canonical lifts. Hence the inverse weight is \(x^{-1}\). Thus the disk-center version of \eqref{eq:M1}--\eqref{eq:M2} gives
\[
                          \operatorname{loc}_wY(\mathbf u)=b^+\zeta
\]
over the generic fields.

To justify transfer here one may first test finite tame character tuples away from the input denominator and polynomial zeros. Nonzero generic minors of the cochain matrices including a putative nonboundary can simultaneously be retained by the bounded-series test on finite characters. At such a fixed tuple the local character field is unramified cyclic \(D_j\), contained locally in the stage base ring fields eventually by the Frobenius exponents. The specialized localization, restricted to \(D_j\), is projected Kummer, whose log evaluated at our embedding is the limit of the full weighted orbit sums. Indeed norm first to \(D_j\) by transitivity grouping the translations; by fixed-group comparison Kummer congruences there give log congruences at increasing precision (fields and projector losses now fixed). These are sums, not averages. Likewise the formal systems for \(\zeta\) evaluate by their weighted logs. This detects the projected finite-character Kummer line: after restriction the log lies in the corresponding character subspace, so the other \(D_j\)-embeddings are translates with unit weights; restriction itself is injective in characteristic zero. Thus the center formula of \eqref{eq:M2} proves the cohomological identity at those tests by all the \(\pi_i\), hence generically. Strict generic acyclicity by \eqref{eq:R1} gives uniqueness from this localization (other non-dyadic allowed terms acyclic there), proving \eqref{eq:R5}.
\end{proof}

\begin{lemma}[Horizontal integrality and Fricke invariance]\label{real:horizontal-integrality}
Before further specialization, the global coordinates of \(z\) are
integral after inverting two. Its all-zero specialization is invariant
under \(w_nc\), componentwise. On a primitive component whose
ordinary characteristic-zero Selmer dimension over \(K\) exceeds
one, this specialization vanishes.
\end{lemma}

\begin{proof}
On any \(f\) and test DVR as in \eqref{eq:R4}, let \(C_H\) use
those dyadic conditions and both unramified conditions at each fixed
odd pair. Moving odd local terms are contractible by their nonzero
component traces, including on the fibers \(u_j=0\); the place
over \(d\) introduces no defect because its Frobenius is a cycle.
The degree-zero injections on fibers and global invariant vanishing
therefore give amplitude \(1,2\). Here invariant vanishing follows
from the evaluation comparison and absolute simplicity over \(K\),
also after the scalar twist. The prescribed local conditions give
exact conjugate self-duality over the DVR.

We first show that generic \(H^1(C_H)\) is a line mapping
isomorphically to the finite line at \(w\). All odd local terms are
generically acyclic. With both dyadic places full, global duality
makes the localization image self-annihilating for the symmetric
conjugate pairing. Write \(V_w,V_{\bar w}\) for the full generic
local cohomology planes. Strict acyclicity identifies the image with
the graph of a map \(A_w:V_w\to V_{\bar w}\), with zero global
localization kernel. Write \(b\) for the cross pairing of these
local planes. For a vector \(x\) in the finite line at \(w\),
isotropy gives
\[
 0=\langle(x,A_wx),(x,A_wx)\rangle=2b(x,A_wx).
\]
Two is invertible in this horizontal test, so \(A_wx\) lies in
the annihilator of that line, namely the chosen opposite finite
line. This proves the assertion about \(H^1(C_H)\). The same
argument applies on the line germ when the generic ranks and
nonconstant fixed-place weights are retained.

Write \(p_0\) for a generator of \(H^1(C_H)\), \(\tau\) for the degree-two torsion length, so \(d(C_H,(p_0,p_0^\vee))=-\tau\) by exact duality. With the same odd unramified conditions but strict/full dyadically, the strict determinant valuation is
\[
                2v(\mathcal P_f)-2\sum_{q\in S^\circ}v(P_q^f(\xi_q)).
\]
Indeed pass from full at the odd places of \(S\) as in \eqref{eq:B1} (both singular factors), using the ratio \((s_F)_f/\mathcal P_f^2\) proved at \eqref{eq:R3}. Now \eqref{eq:D} and \eqref{eq:R4} give
\[
 v(\mathcal P_f)=\tau/2+v(\log\operatorname{loc}_w p_0)
                             +\sum_{q\in S}v(P_q^f(\xi_q)).
\]
In particular the generic \(z_i=\pi_i z\) is the unrestricted global image of a multiple of \(p_0\) of exponent \(\tau/2\), by the local coordinate formula and uniqueness. It is integral here. With no degree-zero unrestricted global term, this gives coordinate integrality after inverting 2 by normality on every integer component; the order and multiplicities split after this inversion with constants sufficiently extended.

Take also a group-power line through zero retaining all the indicated generic conditions by general integer slopes, and apply the same calculation at its characteristic-zero center DVR. The special degree one now gives the usual Kummer Selmer on \(f\) over \(K\), by \eqref{eq:R4}, odd characteristic-zero acyclicity there and unramified inflation. Thus \(z_i(0)\) lies in that space, and is zero if its rank exceeds one (\(\tau>0\)). In the rank-one case complex conjugation on that Selmer line is the primitive Fricke sign. Indeed the two ranks over \(\mathbb Q\) for \(f,f\otimes\chi_{-d}\) add by quadratic induction, and the member with odd functional sign has odd rank by the form parity above; the Fricke sign equals minus the functional sign of \(f\). Moreover \(z_i=\xi_{\ell_0}^i z_0\) on the generic primitive coordinates by the same uniqueness. At old level \(w_n\) exchanges pushforwards \(i,2-i\) with precisely the primitive Fricke action (commuting degeneracy and Fricke on the analytic arguments). Thus \(z(0)\) is \(w_nc\)-invariant, componentwise.
\end{proof}

The last passage to integral coefficients uses two different kinds
of denominator control. We record the elementary intersection that
combines them.

\begin{lemma}[Removing the last denominator]\label{real:denominator-intersection}
Let \(\mathcal O\) be a discrete valuation ring of residue
characteristic two, with uniformizer \(\varpi\), and let
\(B=\mathcal O[[u]]\). Let \(N\) be a finite free \(B\)-module.
If \(v\in N[1/\varpi]\), \(cv\in N\), and the image of
\(c\in B\) in \((\mathcal O/\varpi)[[u]]\) is nonzero, then
\(v\in N\).
\end{lemma}

\begin{proof}
Write \(v=\varpi^{-a}v_0\) in coordinates. The condition is
\(cv_0\in\varpi^aN\). Multiplication by the nonzero residual
series \(\overline c\) is injective on each coordinate of
\(N/\varpi N\), so \(v_0\in\varpi N\). Repeating the argument
\(a\) times gives \(v_0\in\varpi^aN\), as required. Equivalently,
\(B[1/\varpi]\cap B[1/c]=B\) in its fraction field.
\end{proof}

\begin{proposition}[Primitivity detected at infinity]\label{real:primitivity}
The paired elliptic quotient \(V_*\) of \eqref{eq:R3} is an integral
unit.
\end{proposition}

\begin{proof}
Specialize first \(u_j=0\) for \(j>0\), and put
\[
 z_0=z|_{u_j=0\ (j>0)},\qquad
 c_1=c|_{u_j=0\ (j>0)},\qquad Q=\prod_j\alpha_j.
\]
The factor \(\alpha_0\) is a unit. Let \(N\) be the module of
global coordinates over the remaining scalar series ring in
\(u_0\), including the multiplicity lattice. It is free over that
scalar ring. The two preceding constructions give
\[
          z_0/Q\in N[1/2],\qquad c_1z_0/Q\in N.
\]
The first inclusion is horizontal integrality and the nonvanishing
of every constant trace. The second is the cofactor divisibility
\eqref{eq:R2} applied after clearing the input. By
Lemma~\ref{real:input-clearing}, \(\overline c_1\ne0\) on the
residual \(u_0\)-line. Lemma~\ref{real:denominator-intersection}
therefore gives \(z_0/Q\in N\).

Its value at \(u_0=0\) is an integral global class \(p\) over
the extended integers. Lemma~\ref{real:center-inflation} identifies
it with a class at fixed support. It is \(w_nc\)-invariant by
Lemma~\ref{real:horizontal-integrality}; the action comparison uses
the same inflation identification in characteristic zero. The
integral global \(H^1\) is torsion-free, with no degree-zero term,
so this invariance also holds integrally.

At this center \eqref{eq:R5} consequently gives
\[
                                 Y=R_b(0)\,p
\]
integrally in Kummer cohomology on the packet at fixed support. Indeed the added inert traces contribute exactly \(\prod_j\alpha_j\) on \(M\) by the good Hecke correspondences (\(T_{q_j}\infty=(q_j+1)\infty\)), taking the limits of actual \(T_{q_j}\) and unramified inflation. Cancel rationally then use torsion-freeness.

Invariant classes descend integrally to the twist of \(M\) by \(w_n\) along \(K/\mathbb Q\), by inflation--restriction (\(M^{G_K}=0\)); one can use unrestricted absolute groups for this descent. Flat extended coefficients commute with invariants and cohomology on the fixed-support models, so one descends there before scalar combinations. Apply \(T_\ell-\ell-1\) for the detector prime to the equality and restrict the descended classes to infinity. On the left this gives the projected \([Y']\), since the invariant actual point \(Y'\) itself has the twisted Kummer and descent is unique. But \(\Theta\) of \eqref{eq:R0} sees that projection nontrivially and is Hecke-linear via the residue map. A nonunit \(V_*\) would put \(R_b(0)\) in the extended maximal ideal by \eqref{eq:R3}, giving zero detector value on the right. This contradicts \eqref{eq:R0}; therefore \(V_*\) is a unit.
\end{proof}

\subsection{The simple zero and the exact product factor}

\begin{proof}[Proof of Proposition~\ref{real:anchor}]
Use this unit equality first on a tame line at \(t=0\) retaining generic strict acyclicity and nonconstant dyadic exponents, at the characteristic-zero center DVR. The rank-detection argument in the proof of Proposition~\ref{split-local:center-rank} applies, with the following local checks: all other finite allowed places are locally contractible at that germ including the new ones by nonzero elliptic traces. The dyadic comparisons use the unramified Kummer/log version near the center (valid for general reduction there). The disk-center identity at \eqref{eq:M2} gives the weighted Hilbert/ring trace localization square up to units (all Euler factors nonzero at the center). Thus ordinary dyadic finite conditions give generic paired lines by \eqref{eq:D}, with zero class-functional valuation by the unit quotient. Their special ranks are also one by \(s_2(E_0/K)=1\) and inflation. Hence the class specializes nontrivially, i.e. the conductor-one Hilbert trace on \(E_0\) is nontorsion (again times all the extra inert traces, nonzero). This yields a simple product zero by \eqref{eq:GZ-E}.

Finally evaluation of the unit equality at the all-zero center integrally cancels exactly the squared trace factors: for each added inert place changing full to unramified there removes the degree-\((1,2)\) singular block of determinant valuation \(2v_2(a_{j,*})\), \(a_{j,*}\) the corresponding elliptic limiting trace, by Frobenius square minus one (norm tending to 1). On the logs in each orientation the trace at that conductor gives multiplier \(a_{j,*}\). Thus for the true strict complex versus paired formula \eqref{eq:M3} we have precisely the equality used at \eqref{eq:E}, now for \(E_0\) with support \(S\) (the Haar calculation includes extra split good primes identically). Acyclicity over characteristic zero at this center follows as there since the point spans with nonzero dyadic log. Hence \eqref{eq:E} and \eqref{eq:G} give \(X(E_0)+X(E_0^{-d})=0\), proving the split-pair anchor here with the discriminants chosen above. This uses the split Gross--Zagier conventions at the outset.
\end{proof}
